% concatenated from main.tex
%********************************************%
%*       Generated from PreTeXt source      *%
%*       on 2025-12-18T22:15:23-05:00       *%
%*   A recent stable commit (2022-07-01):   *%
%* 6c761d3dba23af92cba35001c852aac04ae99a5f *%
%*                                          *%
%*         https://pretextbook.org          *%
%*                                          *%
%********************************************%
\documentclass[oneside,12pt,]{article}
%% Custom Preamble Entries, early (use latex.preamble.early)
%% Default LaTeX packages
%%   1.  always employed (or nearly so) for some purpose, or
%%   2.  a stylewriter may assume their presence
\usepackage{geometry}
%% Some aspects of the preamble are conditional,
%% the LaTeX engine is one such determinant
\usepackage{ifthen}
%% etoolbox has a variety of modern conveniences
\usepackage{etoolbox}
\usepackage{ifxetex,ifluatex}
%% Raster graphics inclusion
\usepackage{graphicx}
%% Color support, xcolor package
%% Always loaded, for: add/delete text, author tools
%% Here, since tcolorbox loads tikz, and tikz loads xcolor
\PassOptionsToPackage{dvipsnames,svgnames,table}{xcolor}
\usepackage{xcolor}
%% begin: defined colors, via xcolor package, for styling
%% end: defined colors, via xcolor package, for styling
%% Colored boxes, and much more, though mostly styling
%% skins library provides "enhanced" skin, employing tikzpicture
%% boxes may be configured as "breakable" or "unbreakable"
%% "raster" controls grids of boxes, aka side-by-side
\usepackage{tcolorbox}
\tcbuselibrary{skins}
\tcbuselibrary{breakable}
\tcbuselibrary{raster}
%% We load some "stock" tcolorbox styles that we use a lot
%% Placement here is provisional, there will be some color work also
%% First, black on white, no border, transparent, but no assumption about titles
\tcbset{ bwminimalstyle/.style={size=minimal, boxrule=-0.3pt, frame empty,
colback=white, colbacktitle=white, coltitle=black, opacityfill=0.0} }
%% Second, bold title, run-in to text/paragraph/heading
%% Space afterwards will be controlled by environment,
%% independent of constructions of the tcb title
%% Places \blocktitlefont onto many block titles
\tcbset{ runintitlestyle/.style={fonttitle=\blocktitlefont\upshape\bfseries, attach title to upper} }
%% Spacing prior to each exercise, anywhere
\tcbset{ exercisespacingstyle/.style={before skip={1.5ex plus 0.5ex}} }
%% Spacing prior to each block
\tcbset{ blockspacingstyle/.style={before skip={2.0ex plus 0.5ex}} }
%% xparse allows the construction of more robust commands,
%% this is a necessity for isolating styling and behavior
%% The tcolorbox library of the same name loads the base library
\tcbuselibrary{xparse}
%% The tcolorbox library loads TikZ, its calc package is generally useful,
%% and is necessary for some smaller documents that use partial tcolor boxes
%% See:  https://github.com/PreTeXtBook/pretext/issues/1624
\usetikzlibrary{calc}
%% We use some more exotic tcolorbox keys to restore indentation to parboxes
\tcbuselibrary{hooks}
%% Save default paragraph indentation and parskip for use later, when adjusting parboxes
\newlength{\normalparindent}
\newlength{\normalparskip}
\AtBeginDocument{\setlength{\normalparindent}{\parindent}}
\AtBeginDocument{\setlength{\normalparskip}{\parskip}}
\newcommand{\setparstyle}{\setlength{\parindent}{\normalparindent}\setlength{\parskip}{\normalparskip}}%% Hyperref should be here, but likes to be loaded late
%%
%% Inline math delimiters, \(, \), need to be robust
%% 2016-01-31:  latexrelease.sty  supersedes  fixltx2e.sty
%% If  latexrelease.sty  exists, bugfix is in kernel
%% If not, bugfix is in  fixltx2e.sty
%% See:  https://tug.org/TUGboat/tb36-3/tb114ltnews22.pdf
%% and read "Fewer fragile commands" in distribution's  latexchanges.pdf
\IfFileExists{latexrelease.sty}{}{\usepackage{fixltx2e}}
%% Text height identically 9 inches, text width varies on point size
%% See Bringhurst 2.1.1 on measure for recommendations
%% 75 characters per line (count spaces, punctuation) is target
%% which is the upper limit of Bringhurst's recommendations
\geometry{letterpaper,total={408pt,9.0in}}
%% Custom Page Layout Adjustments (use publisher page-geometry entry)
%% This LaTeX file may be compiled with pdflatex, xelatex, or lualatex executables
%% LuaTeX is not explicitly supported, but we do accept additions from knowledgeable users
%% The conditional below provides  pdflatex  specific configuration last
%% begin: engine-specific capabilities
\ifthenelse{\boolean{xetex} \or \boolean{luatex}}{%
%% begin: xelatex and lualatex-specific default configuration
\ifxetex\usepackage{xltxtra}\fi
%% realscripts is the only part of xltxtra relevant to lualatex 
\ifluatex\usepackage{realscripts}\fi
%% end:   xelatex and lualatex-specific default configuration
}{
%% begin: pdflatex-specific default configuration
%% We assume a PreTeXt XML source file may have Unicode characters
%% and so we ask LaTeX to parse a UTF-8 encoded file
%% This may work well for accented characters in Western language,
%% but not with Greek, Asian languages, etc.
%% When this is not good enough, switch to the  xelatex  engine
%% where Unicode is better supported (encouraged, even)
\usepackage[utf8]{inputenc}
%% end: pdflatex-specific default configuration
}
%% end:   engine-specific capabilities
%%
%% Fonts.  Conditional on LaTex engine employed.
%% Default Text Font: The Latin Modern fonts are
%% "enhanced versions of the [original TeX] Computer Modern fonts."
%% We use them as the default text font for PreTeXt output.
%% Automatic Font Control
%% Portions of a document, are, or may, be affected by defined commands
%% These are perhaps more flexible when using  xelatex  rather than  pdflatex
%% The following definitions are meant to be re-defined in a style, using \renewcommand
%% They are scoped when employed (in a TeX group), and so should not be defined with an argument
\newcommand{\divisionfont}{\relax}
\newcommand{\blocktitlefont}{\relax}
\newcommand{\contentsfont}{\relax}
\newcommand{\pagefont}{\relax}
\newcommand{\tabularfont}{\relax}
\newcommand{\xreffont}{\relax}
\newcommand{\titlepagefont}{\relax}
%%
\ifthenelse{\boolean{xetex} \or \boolean{luatex}}{%
%% begin: font setup and configuration for use with xelatex
%% Generally, xelatex is necessary for non-Western fonts
%% fontspec package provides extensive control of system fonts,
%% meaning *.otf (OpenType), and apparently *.ttf (TrueType)
%% that live *outside* your TeX/MF tree, and are controlled by your *system*
%% (it is possible that a TeX distribution will place fonts in a system location)
%%
%% The fontspec package is the best vehicle for using different fonts in  xelatex
%% So we load it always, no matter what a publisher or style might want
%%
\usepackage{fontspec}
%%
%% begin: xelatex main font ("font-xelatex-main" template)
%% Latin Modern Roman is the default font for xelatex and so is loaded with a TU encoding
%% *in the format* so we can't touch it, only perhaps adjust it later
%% in one of two ways (then known by NFSS names such as "lmr")
%% (1) via NFSS with font family names such as "lmr" and "lmss"
%% (2) via fontspec with commands like \setmainfont{Latin Modern Roman}
%% The latter requires the font to be known at the system-level by its font name,
%% but will give access to OTF font features through optional arguments
%% https://tex.stackexchange.com/questions/470008/
%% where-and-how-does-fontspec-sty-specify-the-default-font-latin-modern-roman
%% http://tex.stackexchange.com/questions/115321
%% /how-to-optimize-latin-modern-font-with-xelatex
%%
%% end:   xelatex main font ("font-xelatex-main" template)
%% begin: xelatex mono font ("font-xelatex-mono" template)
%% (conditional on non-trivial uses being present in source)
%% end:   xelatex mono font ("font-xelatex-mono" template)
%% begin: xelatex font adjustments ("font-xelatex-style" template)
%% end:   xelatex font adjustments ("font-xelatex-style" template)
%%
%% Extensive support for other languages
\usepackage{polyglossia}
%% Set main/default language based on pretext/@xml:lang value
%% document language code is "en-US", US English
%% usmax variant has extra hypenation
\setmainlanguage[variant=usmax]{english}
%% Enable secondary languages based on discovery of @xml:lang values
%% Enable fonts/scripts based on discovery of @xml:lang values
%% Western languages should be ably covered by Latin Modern Roman
%% end:   font setup and configuration for use with xelatex
}{%
%% begin: font setup and configuration for use with pdflatex
%% begin: pdflatex main font ("font-pdflatex-main" template)
\usepackage{lmodern}
\usepackage[T1]{fontenc}
%% end:   pdflatex main font ("font-pdflatex-main" template)
%% begin: pdflatex mono font ("font-pdflatex-mono" template)
%% (conditional on non-trivial uses being present in source)
%% end:   pdflatex mono font ("font-pdflatex-mono" template)
%% begin: pdflatex font adjustments ("font-pdflatex-style" template)
%% end:   pdflatex font adjustments ("font-pdflatex-style" template)
%% end:   font setup and configuration for use with pdflatex
}
%% Micromanage spacing, etc.  The named "microtype-options"
%% template may be employed to fine-tune package behavior
\usepackage{microtype}
%% Symbols, align environment, commutative diagrams, bracket-matrix
\usepackage{amsmath}
\usepackage{amscd}
\usepackage{amssymb}
%% allow page breaks within display mathematics anywhere
%% level 4 is maximally permissive
%% this is exactly the opposite of AMSmath package philosophy
%% there are per-display, and per-equation options to control this
%% split, aligned, gathered, and alignedat are not affected
\allowdisplaybreaks[4]
%% allow more columns to a matrix
%% can make this even bigger by overriding with  latex.preamble.late  processing option
\setcounter{MaxMatrixCols}{30}
%%
%%
%% Division Titles, and Page Headers/Footers
%% titlesec package, loading "titleps" package cooperatively
%% See code comments about the necessity and purpose of "explicit" option.
%% The "newparttoc" option causes a consistent entry for parts in the ToC 
%% file, but it is only effective if there is a \titleformat for \part.
%% "pagestyles" loads the  titleps  package cooperatively.
\usepackage[explicit, newparttoc, pagestyles]{titlesec}
%% The companion titletoc package for the ToC.
\usepackage{titletoc}
%% begin: customizations of page styles via the modal "titleps-style" template
%% Designed to use commands from the LaTeX "titleps" package
\pagestyle{plain}
%% end: customizations of page styles via the modal "titleps-style" template
%%
%% Create globally-available macros to be provided for style writers
%% These are redefined for each occurence of each division
\newcommand{\divisionnameptx}{\relax}%
\newcommand{\titleptx}{\relax}%
\newcommand{\subtitleptx}{\relax}%
\newcommand{\shortitleptx}{\relax}%
\newcommand{\authorsptx}{\relax}%
\newcommand{\epigraphptx}{\relax}%
%% Create environments for possible occurences of each division
%% Environment for a PTX "section" at the level of a LaTeX "section"
\NewDocumentEnvironment{sectionptx}{mmmmmmm}
{%
\renewcommand{\divisionnameptx}{#1}%
\renewcommand{\titleptx}{#2}%
\renewcommand{\subtitleptx}{#3}%
\renewcommand{\shortitleptx}{#4}%
\renewcommand{\authorsptx}{#5}%
\renewcommand{\epigraphptx}{#6}%
\section[{#4}]{#2}%
\label{#7}%
}{}%
%% Environment for a PTX "subsection" at the level of a LaTeX "subsection"
\NewDocumentEnvironment{subsectionptx}{mmmmmmm}
{%
\renewcommand{\divisionnameptx}{#1}%
\renewcommand{\titleptx}{#2}%
\renewcommand{\subtitleptx}{#3}%
\renewcommand{\shortitleptx}{#4}%
\renewcommand{\authorsptx}{#5}%
\renewcommand{\epigraphptx}{#6}%
\subsection[{#4}]{#2}%
\label{#7}%
}{}%
%% Environment for a PTX "references" at the level of a LaTeX "section"
\NewDocumentEnvironment{references-section}{mmmmmmm}
{%
\renewcommand{\divisionnameptx}{#1}%
\renewcommand{\titleptx}{#2}%
\renewcommand{\subtitleptx}{#3}%
\renewcommand{\shortitleptx}{#4}%
\renewcommand{\authorsptx}{#5}%
\renewcommand{\epigraphptx}{#6}%
\section[{#4}]{#2}%
\label{#7}%
}{}%
%% Environment for a PTX "references" at the level of a LaTeX "section"
\NewDocumentEnvironment{references-section-numberless}{mmmmmmm}
{%
\renewcommand{\divisionnameptx}{#1}%
\renewcommand{\titleptx}{#2}%
\renewcommand{\subtitleptx}{#3}%
\renewcommand{\shortitleptx}{#4}%
\renewcommand{\authorsptx}{#5}%
\renewcommand{\epigraphptx}{#6}%
\section*{#2}%
\addcontentsline{toc}{section}{#4}
\label{#7}%
}{}%
%%
%% Styles for six traditional LaTeX divisions
\titleformat{\part}[display]
{\divisionfont\Huge\bfseries\centering}{\divisionnameptx\space\thepart}{30pt}{\Huge#1}
[{\Large\centering\authorsptx}]
\titleformat{\chapter}[display]
{\divisionfont\huge\bfseries}{\divisionnameptx\space\thechapter}{20pt}{\Huge#1}
[{\Large\authorsptx}]
\titleformat{name=\chapter,numberless}[display]
{\divisionfont\huge\bfseries}{}{0pt}{#1}
[{\Large\authorsptx}]
\titlespacing*{\chapter}{0pt}{50pt}{40pt}
\titleformat{\section}[hang]
{\divisionfont\Large\bfseries}{\thesection}{1ex}{#1}
[{\large\authorsptx}]
\titleformat{name=\section,numberless}[block]
{\divisionfont\Large\bfseries}{}{0pt}{#1}
[{\large\authorsptx}]
\titlespacing*{\section}{0pt}{3.5ex plus 1ex minus .2ex}{2.3ex plus .2ex}
\titleformat{\subsection}[hang]
{\divisionfont\large\bfseries}{\thesubsection}{1ex}{#1}
[{\normalsize\authorsptx}]
\titleformat{name=\subsection,numberless}[block]
{\divisionfont\large\bfseries}{}{0pt}{#1}
[{\normalsize\authorsptx}]
\titlespacing*{\subsection}{0pt}{3.25ex plus 1ex minus .2ex}{1.5ex plus .2ex}
\titleformat{\subsubsection}[hang]
{\divisionfont\normalsize\bfseries}{\thesubsubsection}{1em}{#1}
[{\small\authorsptx}]
\titleformat{name=\subsubsection,numberless}[block]
{\divisionfont\normalsize\bfseries}{}{0pt}{#1}
[{\normalsize\authorsptx}]
\titlespacing*{\subsubsection}{0pt}{3.25ex plus 1ex minus .2ex}{1.5ex plus .2ex}
\titleformat{\paragraph}[hang]
{\divisionfont\normalsize\bfseries}{\theparagraph}{1em}{#1}
[{\small\authorsptx}]
\titleformat{name=\paragraph,numberless}[block]
{\divisionfont\normalsize\bfseries}{}{0pt}{#1}
[{\normalsize\authorsptx}]
\titlespacing*{\paragraph}{0pt}{3.25ex plus 1ex minus .2ex}{1.5em}
%%
%% Styles for five traditional LaTeX divisions
\titlecontents{part}%
[0pt]{\contentsmargin{0em}\addvspace{1pc}\contentsfont\bfseries}%
{\Large\thecontentslabel\enspace}{\Large}%
{}%
[\addvspace{.5pc}]%
\titlecontents{chapter}%
[0pt]{\contentsmargin{0em}\addvspace{1pc}\contentsfont\bfseries}%
{\large\thecontentslabel\enspace}{\large}%
{\hfill\bfseries\thecontentspage}%
[\addvspace{.5pc}]%
\dottedcontents{section}[3.8em]{\contentsfont}{2.3em}{1pc}%
\dottedcontents{subsection}[6.1em]{\contentsfont}{3.2em}{1pc}%
\dottedcontents{subsubsection}[9.3em]{\contentsfont}{4.3em}{1pc}%
%%
%% Begin: Semantic Macros
%% To preserve meaning in a LaTeX file
%%
%% \mono macro for content of "c", "cd", "tag", etc elements
%% Also used automatically in other constructions
%% Simply an alias for \texttt
%% Always defined, even if there is no need, or if a specific tt font is not loaded
\newcommand{\mono}[1]{\texttt{#1}}
%%
%% Following semantic macros are only defined here if their
%% use is required only in this specific document
%%
%% Used for inline definitions of terms
\newcommand{\terminology}[1]{\textbf{#1}}
%% Titles of longer works (e.g. books, versus articles)
\newcommand{\pubtitle}[1]{\textsl{#1}}
%% End: Semantic Macros
%% Equation Numbering
%% Controlled by  numbering.equations.level  processing parameter
%% No adjustment here implies document-wide numbering
\numberwithin{equation}{section}
%% More flexible list management, esp. for references
%% But also for specifying labels (i.e. custom order) on nested lists
\usepackage{enumitem}
%% Lists of references in their own section, maximum depth 1
\newlist{referencelist}{description}{4}
\setlist[referencelist]{leftmargin=!,labelwidth=!,labelsep=0ex,itemsep=1.0ex,topsep=1.0ex,partopsep=0pt,parsep=0pt}
%% hyperref driver does not need to be specified, it will be detected
%% Footnote marks in tcolorbox have broken linking under
%% hyperref, so it is necessary to turn off all linking
%% It *must* be given as a package option, not with \hypersetup
\usepackage[hyperfootnotes=false]{hyperref}
%% For a print PDF, no surrounding boxes, so simply textcolor (but still active to preserve spacing)
\hypersetup{hidelinks}
%% Less-clever names for hyperlinks are more reliable, *especially* for structural parts
%% See comments in the code to learn more about the importance of this setting
\hypersetup{hypertexnames=false}
%%The  hypertexnames  setting then confuses the hyperlinking from the index
%%This patch resolves the incorrect links, see code for StackExchange post.
\makeatletter
\patchcmd\Hy@EveryPageBoxHook{\Hy@EveryPageAnchor}{\Hy@hypertexnamestrue\Hy@EveryPageAnchor}{}{\fail}
\makeatother
\hypersetup{pdftitle={Correctness of Extended RSA Public Key Cryptosystem}}
%% If you manually remove hyperref, leave in this next command
%% This will allow LaTeX compilation, employing this no-op command
\providecommand\phantomsection{}
%% Division Numbering: Chapters, Sections, Subsections, etc
%% Division numbers may be turned off at some level ("depth")
%% A section *always* has depth 1, contrary to us counting from the document root
%% The latex default is 3.  If a larger number is present here, then
%% removing this command may make some cross-references ambiguous
%% The precursor variable $numbering-maxlevel is checked for consistency in the common XSL file
\setcounter{secnumdepth}{3}
%%
%%
%% A faux tcolorbox whose only purpose is to provide common numbering
%% facilities for most blocks (possibly not projects, 2D displays)
%% Controlled by  numbering.theorems.level  processing parameter
\newtcolorbox[auto counter, number within=section]{block}{}
%%
%% This document is set to number PROJECT-LIKE on a separate numbering scheme
%% So, a faux tcolorbox whose only purpose is to provide this numbering
%% Controlled by  numbering.projects.level  processing parameter
\newtcolorbox[auto counter, number within=section]{project-distinct}{}
%% A faux tcolorbox whose only purpose is to provide common numbering
%% facilities for 2D displays which are subnumbered as part of a "sidebyside"
\makeatletter
\newtcolorbox[auto counter, number within=tcb@cnt@block, number freestyle={\noexpand\thetcb@cnt@block(\noexpand\alph{\tcbcounter})}]{subdisplay}{}
\makeatother
%%
%% tcolorbox, with styles, for THEOREM-LIKE
%%
%% theorem: fairly simple numbered block/structure
\tcbset{ theoremstyle/.style={bwminimalstyle, runintitlestyle, blockspacingstyle, after title={\space}, before upper app={\setparstyle}, } }
\newtcolorbox[use counter from=block]{theorem}[4]{title={{#1~\thetcbcounter\notblank{#2#3}{\space}{}\notblank{#2}{\space#2}{}\notblank{#3}{\space(#3)}{}}}, phantomlabel={#4}, breakable, after={\par}, fontupper=\itshape, theoremstyle, }
%% corollary: fairly simple numbered block/structure
\tcbset{ corollarystyle/.style={bwminimalstyle, runintitlestyle, blockspacingstyle, after title={\space}, before upper app={\setparstyle}, } }
\newtcolorbox[use counter from=block]{corollary}[4]{title={{#1~\thetcbcounter\notblank{#2#3}{\space}{}\notblank{#2}{\space#2}{}\notblank{#3}{\space(#3)}{}}}, phantomlabel={#4}, breakable, after={\par}, fontupper=\itshape, corollarystyle, }
%%
%% tcolorbox, with styles, for PROOF-LIKE
%%
%% proof: title is a replacement
\tcbset{ proofstyle/.style={bwminimalstyle, fonttitle=\blocktitlefont\itshape, attach title to upper, after title={\space}, after upper={\space\space\hspace*{\stretch{1}}\(\blacksquare\)},
} }
\newtcolorbox{proof}[3]{title={\notblank{#2}{#2}{#1.}}, phantom={\label{#3}\hypertarget{#3}{}}, breakable, after={\par}, proofstyle, before upper app={\setparstyle} }
%%
%% tcolorbox, with styles, for DEFINITION-LIKE
%%
%% definition: fairly simple numbered block/structure
\tcbset{ definitionstyle/.style={bwminimalstyle, runintitlestyle, blockspacingstyle, after title={\space}, after upper={\space\space\hspace*{\stretch{1}}\(\lozenge\)}, before upper app={\setparstyle}, } }
\newtcolorbox[use counter from=block]{definition}[3]{title={{#1~\thetcbcounter\notblank{#2}{\space\space#2}{}}}, phantomlabel={#3}, breakable, after={\par}, definitionstyle, }
%%
%% tcolorbox, with styles, for REMARK-LIKE
%%
%% remark: fairly simple numbered block/structure
\tcbset{ remarkstyle/.style={bwminimalstyle, runintitlestyle, blockspacingstyle, after title={\space}, before upper app={\setparstyle}, } }
\newtcolorbox[use counter from=block]{remark}[3]{title={{#1~\thetcbcounter\notblank{#2}{\space\space#2}{}}}, phantomlabel={#3}, breakable, after={\par}, remarkstyle, }
%%
%% tcolorbox, with styles, for EXAMPLE-LIKE
%%
%% example: fairly simple numbered block/structure
\tcbset{ examplestyle/.style={bwminimalstyle, runintitlestyle, blockspacingstyle, after title={\space}, after upper={\space\space\hspace*{\stretch{1}}\(\square\)}, before upper app={\setparstyle}, } }
\newtcolorbox[use counter from=block]{example}[3]{title={{#1~\thetcbcounter\notblank{#2}{\space\space#2}{}}}, phantomlabel={#3}, breakable, after={\par}, examplestyle, }
%%
%% xparse environments for introductions and conclusions of divisions
%%
%% introduction: in a structured division
\NewDocumentEnvironment{introduction}{m}
{\notblank{#1}{\noindent\textbf{#1}\space}{}}{\par\medskip}
%% Graphics Preamble Entries
\usepackage{tikz}
%% If tikz has been loaded, replace ampersand with \amp macro
%% Custom Preamble Entries, late (use latex.preamble.late)
%% extpfeil package for certain extensible arrows,
%% as also provided by MathJax extension of the same name
%% NB: this package loads mtools, which loads calc, which redefines
%%     \setlength, so it can be removed if it seems to be in the 
%%     way and your math does not use:
%%     
%%     \xtwoheadrightarrow, \xtwoheadleftarrow, \xmapsto, \xlongequal, \xtofrom
%%     
%%     we have had to be extra careful with variable thickness
%%     lines in tables, and so also load this package late
\usepackage{extpfeil}
%% Begin: Author-provided macros
%% (From  docinfo/macros  element)
%% Plus three from PTX for XML characters
\newcommand{\foo}{b^{ar}}
\newcommand{\lt}{<}
\newcommand{\gt}{>}
\newcommand{\amp}{&}
%% End: Author-provided macros
%% Title page information for article
\title{Correctness of Extended RSA Public Key Cryptosystem}
\author{Dar-jen Chang\\
University of Louisville\\
\href{mailto:djchan01@louisville.edu}{\nolinkurl{djchan01@louisville.edu}}
\and
Suranjan Gautam\\
University of Louisville\\
\href{mailto:suranjan.gautam@louisville.edu}{\nolinkurl{suranjan.gautam@louisville.edu}}
}
\date{}
\begin{document}
%% bottom alignment is explicit, since it normally depends on oneside, twoside
\raggedbottom
%% Target for xref to top-level element is document start
\label{shorttitlelowercase}\hypertarget{shorttitlelowercase}{}
\maketitle
\thispagestyle{empty}
\renewcommand*{\abstractname}{Abstract}
\begin{abstract}
This paper proposes an alternative approach to formally establishing the correctness of the RSA public key cryptosystem. The methodology presented herein deviates slightly from conventional proofs found in existing literature. Specifically, this study explores the conditions under which the choice of the positive integer N, a fundamental component of RSA, can be extended beyond the standard selection criteria. We derive explicit conditions that determine when certain values of N are valid for the encryption scheme and explain why others may fail to satisfy the correctness requirements. The scope of this paper is limited to the mathematical proof of correctness for RSA-like schemes, deliberately omitting issues related to the cryptographic security of RSA.%
\end{abstract}
%
%
\typeout{************************************************}
\typeout{Section 1 Introduction and Definitions}
\typeout{************************************************}
%
\begin{sectionptx}{Section}{Introduction and Definitions}{}{Introduction and Definitions}{}{}{section-1}
\begin{introduction}{}%
This section introduces the mathematical framework for our extended RSA encryption scheme and establishes the correctness conditions that will be analyzed throughout this paper.%
\end{introduction}%
%
%
\typeout{************************************************}
\typeout{Subsection 1.1 Problem Statement}
\typeout{************************************************}
%
\begin{subsectionptx}{Subsection}{Problem Statement}{}{Problem Statement}{}{}{subsection-1}
\begin{introduction}{}%
To achieve our goals, we begin by defining an extended framework for the RSA public key encryption system.%
\end{introduction}%
\begin{definition}{Definition}{Extended RSA Public Key Encryption Scheme.}{def-generalized-rsa}%
Given a positive integer \(N\) and \(\phi(N)\) denoting the Euler's totient function, we choose an integer \(e\) such that \(1 \leq e < \phi(N)\) and \(\gcd(e, \phi(N)) = 1\). We then compute the multiplicative inverse of \(e\) modulo \(\phi(N)\) and denote it as \(d\), where \(ed \equiv 1 \pmod{\phi(N)}\).%
\par
For any integer \(m\), \(1 \leq m < N\), the encryption and decryption functions are defined as follows:%
\par
\terminology{Encryption Function (\(Enc\)):} \(c = Enc(m) = m^e \pmod N\)%
\par
\terminology{Decryption Function (\(Dec\)):} \(m = Dec(c) = c^d \pmod N\)%
\end{definition}
\begin{remark}{Remark}{}{subsection-1-4}%
The pair \((e, N)\) forms the public key, while \((d, N)\) forms the private key.%
\end{remark}
\begin{definition}{Definition}{Correctness of Generalized RSA Public Key Encryption Scheme.}{def-generalized-rsa-2}%
The correctness condition for this encryption scheme is described as follows:%
\par
For any integer \(m\), \(1 \leq m < N\), we have:%
%
\begin{equation}
Dec(Enc(m)) \equiv m \pmod{N}\label{eq-correctness}
\end{equation}
or equivalently,%
%
\begin{equation}
(m^e)^d \equiv m \pmod{N}\label{def-generalized-rsa-2-2-5}
\end{equation}
\end{definition}
The classical (or textbook) RSA algorithm chooses \(N\) as the product of two large, distinct prime numbers, which guarantees that \hyperref[eq-correctness]{({\xreffont\ref{eq-correctness}})} holds. However, this paper seeks to investigate the generalized criteria under which \hyperref[eq-correctness]{({\xreffont\ref{eq-correctness}})} remains valid for various choices of \(N\). To achieve this, we provide a modified proof of correctness that builds upon existing proofs in the literature and derive necessary criteria for the selection of \(N\).%
\end{subsectionptx}
\end{sectionptx}
%
%
\typeout{************************************************}
\typeout{Section 2 Preliminary Number Theory Background}
\typeout{************************************************}
%
\begin{sectionptx}{Section}{Preliminary Number Theory Background}{}{Preliminary Number Theory Background}{}{}{section-2}
\begin{introduction}{}%
In this section, we will cover some basic concepts of number theory such as divisibility, fundamental theorem of arithmetic, greatest common divisor, and modular arithmetic, which are important for understanding RSA.%
\end{introduction}%
%
%
\typeout{************************************************}
\typeout{Subsection 2.1 Integer Division and Fundamental theorem of Arithmetic}
\typeout{************************************************}
%
\begin{subsectionptx}{Subsection}{Integer Division and Fundamental theorem of Arithmetic}{}{Integer Division and Fundamental theorem of Arithmetic}{}{}{subsection-2-1}
Integer division and remainder are fundamental operations in modular arithmetic. Given two integers \(a\) and \(b\), where \(b \neq 0\), the integer division of \(a\) by \(b\) is denoted as \(q = a \div b\), and the remainder is denoted as \(r = a \bmod b\). These values satisfy the equation:%
\par
\(a = b \cdot q + r\), where \(0 \leq r < b\).%
\par
If the remainder \(r\) is zero, then \(a\) is divisible by \(b\), and we say that \(b\) divides \(a\) and is denoted by \(b|a\).%
\begin{theorem}{Theorem}{Fundamental theorem of Arithmetic.}{}{Factorization-Theorem}%
Every integer \(N\) greater than 1 can be uniquely expressed as a product of prime factors:%
\begin{equation}
N = p_1^{n_1} p_2^{n_2} \cdots p_k^{n_k} = \prod_{i=1}^{k} p_i^{n_i}\label{Factorization-Theorem-2-1-2}
\end{equation}
where \(p_1,p_2,\ldots,p_k\) are distinct prime numbers and \(n_1,n_2,\ldots,n_k\) are their corresponding positive integer exponents.%
\end{theorem}
If%
\begin{equation*}
N = p_1^{n_1} p_2^{n_2} \cdots p_k^{n_k}
\end{equation*}
and \(a|N\), then \(a\) can be expressed as a product of the prime factors:%
\begin{equation*}
a = p_1^{m_1} p_2^{m_2} \cdots p_k^{m_k}
\end{equation*}
where \(0 \leq m_i \leq n_i\) (usually if \(m_i = 0\), \(p_i\) is not included in the product). This means that \(a\) can be formed by taking some of the prime factors of \(N\) and raising them to powers that do not exceed the corresponding powers in the prime factorization of \(N\).%
\end{subsectionptx}
%
%
\typeout{************************************************}
\typeout{Subsection 2.2 The Greatest Common Divisor and Co-Prime}
\typeout{************************************************}
%
\begin{subsectionptx}{Subsection}{The Greatest Common Divisor and Co-Prime}{}{The Greatest Common Divisor and Co-Prime}{}{}{subsection-2-2}
The greatest common divisor (gcd) of two integers \(a\) and \(b\), denoted as \(\gcd(a, b)\), is the largest positive integer that divides both \(a\) and \(b\) without leaving a remainder.%
\par
Two integers \(a\) and \(b\) are said to be co-prime (or relatively prime) if their gcd is 1, i.e., \(\gcd(a, b) = 1\). This implies that \(a\) and \(b\) have no common positive integer factors other than 1.%
\par
Many useful results in number theory and cryptography rely on the properties of co-prime integers. For example, in the RSA public key cryptosystem, the choice of the public exponent \(e\) must be co-prime to the Euler totient function \(\phi(N)\) to ensure that the encryption and decryption processes work correctly. To prove the correctness of the RSA algorithm, we need the following two theorems:%
\begin{theorem}{Theorem}{Product of Divisors.}{}{product-of-divisors}%
For any three integers \(a\), \(b\), and \(N\), if \(a|N\), \(b|N\) and \(\gcd(a, b) = 1\), then \(ab|N\).%
\end{theorem}
\begin{proof}{Proof}{}{product-of-divisors-3}
Assume%
\begin{equation*}
N = p_1^{n_1} p_2^{n_2} \cdots p_k^{n_k}
\end{equation*}
as given in the Fundamental theorem of Arithmetic. Since \(a|N\) and \(b|N\), we can express \(a\) and \(b\) in terms of the prime factors of \(N\):%
\begin{equation*}
a = p_1^{m_1} p_2^{m_2} \cdots p_k^{m_k}
\end{equation*}
and%
\begin{equation*}
b = p_1^{l_1} p_2^{l_2} \cdots p_k^{l_k}
\end{equation*}
where \(0 \leq m_i \leq n_i\) and \(0 \leq l_i \leq n_i\). Since \(\gcd(a, b) = 1\), it follows that for each prime factor \(p_i\), either \(m_i = 0\) or \(l_i = 0\). Therefore, the product \(ab\) can be expressed as:%
\begin{equation*}
ab = p_1^{m_1 + l_1} p_2^{m_2 + l_2} \cdots p_k^{m_k + l_k}
\end{equation*}
where \(m_i + l_i \leq n_i\) for each \(i\). This implies that \(ab\) is a product of the prime factors of \(N\) raised to powers that do not exceed the corresponding powers in the prime factorization of \(N\). Hence, \(ab|N\).%
\end{proof}
As a consequence, the product of two coprime integers divides any common multiple of those integers. This theorem is particularly useful in the context of RSA, where we often deal with products of two coprime integers.%
\begin{theorem}{Theorem}{}{}{co-prime-of-divisors}%
Given any two coprime integers \(P\) and \(Q\) and any two integers \(p\) and \(q\), if \(p|P\) and \(q|Q\), then \(p\) and \(q\) are coprime, i.e., \(\gcd(p, q) = 1\).%
\end{theorem}
\begin{proof}{Proof}{}{co-prime-of-divisors-2}
This theorem can be proved using similar reasoning as in \hyperref[product-of-divisors]{Theorem~{\xreffont\ref{product-of-divisors}}, p.\,\pageref{product-of-divisors}}.%
\end{proof}
\begin{theorem}{Theorem}{}{}{product_of_gcd}%
For any three positive integers \(m\), \(P\), and \(Q\), we have:%
%
\begin{align*}
(1) \quad & \gcd(m, PQ) \leq \gcd(m, P) \times \gcd(m, Q)\\
(2) \quad & \text{ if } \gcd(P, Q) = 1, \text{ then } \gcd(m, PQ) = \gcd(m, P) \times \gcd(m, Q)
\end{align*}
\begin{proof}{Proof}{}{product_of_gcd-1-3}
(1)%
\par
By the definition of gcd, we have:%
\begin{equation*}
g = \gcd(m, PQ) = \max\{d \mid d \text{ divides } m \text{ and } d \text{ divides } PQ\}
\end{equation*}
Since g divides \(PQ\) and \(m\), g can be expressed as \(p * q\), where \(p\) divides \(P\) and \(m\), and \(q\) divides \(Q\) and \(m\). Therefore,%
\begin{equation*}
p \leq \gcd(m, P) \quad \text{and} \quad q \leq \gcd(m, Q)
\end{equation*}
which implies that%
\begin{equation*}
g = p * q \leq \gcd(m, P) * \gcd(m, Q)
\end{equation*}
%
\par
This proves part (1).%
\par
(2)%
\par
Let%
\begin{equation*}
p = \gcd(m, P) \quad \text{and} \quad q = \gcd(m, Q)
\end{equation*}
Since \(\gcd(P, Q) = 1\), by \hyperref[co-prime-of-divisors]{Theorem~{\xreffont\ref{co-prime-of-divisors}}, p.\,\pageref{co-prime-of-divisors}}, we have \(\gcd(p, q) = 1\).%
\par
Consequently, by \hyperref[product-of-divisors]{Theorem~{\xreffont\ref{product-of-divisors}}, p.\,\pageref{product-of-divisors}}, we have \(p*q\) divides both \(m\) and \(PQ\). Therefore,%
\begin{equation*}
\gcd(m,P) \times \gcd(m,Q) = p \times q \leq \gcd(m, PQ) 
\end{equation*}
Together with (1), we prove part (2).%
\end{proof}
\end{theorem}
\end{subsectionptx}
%
%
\typeout{************************************************}
\typeout{Subsection 2.3 Modular Arithmetc and Congruences}
\typeout{************************************************}
%
\begin{subsectionptx}{Subsection}{Modular Arithmetc and Congruences}{}{Modular Arithmetc and Congruences}{}{}{subsection-2-3}
Modular arithmetic is a system of arithmetic for integers, where numbers "wrap around" after reaching a certain value, known as the modulus. More specifically, given a positive integer \(N\), we can partition integers into \(N\) equivalence classes denoted by \({0},{1}, \cdots, {N-1}\)  so that each \({i}\) represents the set of integers that have the same remainder \(i\) when they are divided by \(N\).%
\begin{definition}{Definition}{Congruence Relation.}{def-congruence}%
We say that two integers \(a\) and \(b\) are congruent modulo \(N\), denoted as \(a \equiv b \pmod{N}\), if they have the same remainder when divided by \(N\). Equivalently, \(a \equiv b \pmod{N}\) if and only if \(N | (a - b)\).%
\end{definition}
\begin{theorem}{Theorem}{}{}{congruence-theorem}%
We are given a positive integer \(N\) and%
\begin{equation*}
N = P \times Q
\end{equation*}
where \(P\) and \(Q\) are positive integers and \(\gcd(P, Q) = 1\). The following property holds for any two integers \(a\) and \(b\):%
\begin{equation*}
a \equiv b \pmod{N}
\end{equation*}
if and only if%
\begin{equation*}
a \equiv b \pmod{P} \quad \text{and} \quad a \equiv b \pmod{Q}
\end{equation*}
%
\end{theorem}
\begin{proof}{Proof}{}{subsection-2-3-5}
Based on the fact that%
\begin{equation*}
a \equiv b \pmod{N} \text{ if and only if } N|(a-b)
\end{equation*}
and \hyperref[product-of-divisors]{Theorem~{\xreffont\ref{product-of-divisors}}, p.\,\pageref{product-of-divisors}}, we can easily prove the theorem. \(\square\)%
\end{proof}
This theorem is particularly useful in RSA, where we often work with products of primes and their multiples. It allows us to break down congruences modulo a composite number into simpler congruences modulo its prime factors. This fact is essntial for the proof of correctness of RSA, as it enables us to analyze the encryption and decryption processes separately for each prime factor of the modulus \(N\).%
\end{subsectionptx}
\end{sectionptx}
%
%
\typeout{************************************************}
\typeout{Section 3 \(\phi\text{-set}\) and \(\phi\) (Euler's Totient) Function}
\typeout{************************************************}
%
\begin{sectionptx}{Section}{\(\phi\text{-set}\) and \(\phi\) (Euler's Totient) Function}{}{\(\phi\text{-set}\) and \(\phi\) (Euler's Totient) Function}{}{}{section-3}
\begin{introduction}{}%
In this section, we introduce Euler's Totient Function, which is an important concept in number theory and plays a significant role in the RSA encryption scheme.%
\end{introduction}%
\begin{definition}{Definition}{\(\phi\text{-set}\).}{def-phi-set}%
The \(\phi\text{-set}\) of a positive integer \(N\), denoted as \(\phi\text{-set}(N)\), is the set of all integers from 1 to \(N\) that are coprime to \(N\).%
\end{definition}
For example, if \(N = 10\), then \(\phi\text{-set}(10)\) = \textbraceleft{}1, 3, 7, 9\textbraceright{}, since these are exactly the integers between 1 and 10 that are coprime to 10. In the literature, the \(\phi\text{-set}(N)\) is commonly called the reduced residue system modulo \(N\). It can be verified \(\phi\text{-set}(N)\) forms a group under multiplication modulo \(N\). We use the notation \(\phi\text{-set}(N)\) because we will later extend this set to a bigger set called \(\Phi\text{-set}(N)\), which plays an essential role in the correctness of the RSA encryption scheme.%
\begin{definition}{Definition}{Euler's Totient Function.}{def-euler-totient}%
Euler's Totient Function of a positive integer \(N\), denoted as \(\phi(N)\), is defined as the number of integers from 1 to \(N\) that are coprime to \(N\). As a result \(\phi(N)\) is the number of elements in the \(\phi\text{-set}(N)\).%
\end{definition}
\begin{theorem}{Theorem}{Computing Euler's Totient Function.}{}{Euler-Totient-Theorem}%
%
\begin{enumerate}
\item{}If \(p\) is a prime number, then \(\phi(p) = p - 1\).%
\item{}If \(p_1, p_2, \ldots, p_k\) are distinct prime numbers, then%
\begin{equation*}
\phi(p_1 p_2 \cdots p_k) = (p_1 - 1)(p_2 - 1) \cdots (p_k - 1)\text{.}
\end{equation*}
%
\item{}If \(n\) is a positive integer with the prime factorization \(n = p_1^{k_1} p_2^{k_2} \cdots p_m^{k_m}\), then%
\begin{equation*}
\phi(n) = n \prod_{i=1}^{m} \left( 1 - \frac{1}{p_i} \right)\text{.}
\end{equation*}
%
\end{enumerate}
\end{theorem}
\begin{corollary}{Corollary}{}{}{corollary-euler-totient}%
If \(N = P \times Q\), where gcd(P, Q) = 1, then \(\phi(N) = \phi(P) \cdot \phi(Q)\).%
\end{corollary}
According to the above theorem, we can design an algorithm to compute the \(\phi(N)\) function by using the prime factorization of \(N\). Consequently, the problem of computing \(\phi(N)\) is reduced to the challenge of finding the prime factorization of \(N\). The security of the RSA encryption scheme relies on the computational difficulty of this factorization problem..%
\end{sectionptx}
%
%
\typeout{************************************************}
\typeout{Section 4 \(\Phi\text{-set}\) and \(\Phi\) Function}
\typeout{************************************************}
%
\begin{sectionptx}{Section}{\(\Phi\text{-set}\) and \(\Phi\) Function}{}{\(\Phi\text{-set}\) and \(\Phi\) Function}{}{}{section-4}
\begin{introduction}{}%
In this section, we define the \(\Phi\text{-set}\) of an positive integer \(N\), which is an extension or a super set of \(\phi\text{-set}\) of \(N\) and plays an essential role in the correctness of the RSA encryption scheme.%
\end{introduction}%
\begin{definition}{Definition}{\(\Phi\text{-set}\).}{def-big-phi-set}%
The \(\Phi\text{-set}\) of a positive integer \(N\), denoted as \(\Phi\text{-set}(N)\), is the set of all integers \(m, 1 \leq m \leq N\) that satisfy the following condition:%
\par
%
\begin{equation*}
\text{if } P = \gcd(m, N) \text{ and } Q = N \div P \text{ then } \gcd(P, Q) = 1\text{.}
\end{equation*}
Or equivalently,%
\begin{equation*}
\Phi\text{-set}(N) = \{m, 1 \leq m \leq N \mid \gcd(P,Q) = 1, \text{where } P = \gcd(m, N) \text{ and } Q = N \div P\}
\end{equation*}
%
\end{definition}
For example, if \(N = 10\), then \(\Phi\text{-set}(10)\) = \textbraceleft{}1,2,3,4,5,6,7,8,9,10\textbraceright{}. We can easily verify the result manually. Or we can use the theoream \hyperref[Phi-Set-Theorem]{Theorem~{\xreffont\ref{Phi-Set-Theorem}}, p.\,\pageref{Phi-Set-Theorem}} for \(10 = 2 \ldots 5\) to prove the result.%
\begin{remark}{Remark}{}{section-4-5}%
For \(m \in \Phi\text{-set}(N)\), from the definition and \hyperref[product_of_gcd]{Theorem~{\xreffont\ref{product_of_gcd}}, p.\,\pageref{product_of_gcd}}, it is clear that \(gcd(m, Q) = 1\).%
\end{remark}
The \(\Phi\) function of a positive integer \(N\), denoted as \(\Phi(N)\), is defined as the number of integers in \(\Phi\text{-set}(N)\).%
\par
It is clear that \(\Phi\text{-set}(N)\) is a super set of \(\phi\text{-set}(N)\), i.e., \(\Phi\text{-set}(N) \supseteq \phi\text{-set}(N)\). As a result, we have \(\Phi(N) \geq \phi(N)\).%
\begin{theorem}{Theorem}{Computing \(\Phi\) Function.}{}{Phi-Set-Theorem}%
For an positive integer \(N\) if :%
\begin{equation*}
N = p_1 p_2 \cdots p_k
\end{equation*}
where \(p_1,p_2,\ldots,p_k\) are distinct prime numbers, then \(\Phi(N) = N\), i.e. the \(\Phi\text{-set}(N)\) contains all integers from 1 to \(N\).%
\end{theorem}
\begin{proof}{Proof}{}{Phi-Set-Theorem-3}
Let \(m\) be any integer such that \(1 \leq m \leq N\). We need to show that \(m \in \Phi\text{-set}(N)\). Let \(P = \gcd(m, N)\) and \(Q = N \div P\). Since \(N\) is a product of distinct primes and \(P\) is a divisor of \(N\), \(P\) must be a product of a subset of the prime factors of \(N\). Therefore, \(Q = N \div P\) is the product of the remaining prime factors of \(N\). Cosequently, \(P\) and \(Q\) do not share any common prime factors, which implies that \(\gcd(P, Q) = 1\) and as a result \(m \in \Phi\text{-set}(N)\). This concludes the proof.%
\end{proof}
\end{sectionptx}
%
%
\typeout{************************************************}
\typeout{Section 5 Euler Theorem}
\typeout{************************************************}
%
\begin{sectionptx}{Section}{Euler Theorem}{}{Euler Theorem}{}{}{section-5}
\begin{introduction}{}%
In this section, we state the well-known Euler’s Theorem, a fundamental result in number theory related to modular arithmetic. This theorem is essential for understanding the RSA encryption scheme and its generalizations.%
\end{introduction}%
\begin{theorem}{Theorem}{Euler's Theorem.}{}{Euler-Theorem}%
For any two positive integers \(m \text{ and } N\), if \(\gcd(m,N) = 1\), then \(m^{\phi(n)} \equiv 1 \pmod N\).%
\end{theorem}
\begin{proof}{Proof}{}{Euler-Theorem-3}
The proof of Euler's Theorem is equvalant to prove for any \(m \in \phi\text{-set}(N)\),%
\begin{equation*}
m^{\phi(n)} \equiv 1 \pmod N\text{.}
\end{equation*}
%
\par
Since \(\phi\text{-set}(N)\) forms a group under multiplication modulo \(N\), for any \(m \in \phi\text{-set}(N)\), the cyclic group generated by \(m\) is a subgroup of \(\phi\text{-set}(N)\). Therefore, the order of \(m\), denoted by \(order(m)\), is a divisor of \(\phi(N)\), which is the order of the group \(\phi\text{-set}(N)\).%
\par
Therefore, we have%
\par
%
\begin{equation*}
m^{order(m)} \equiv 1 \pmod N
\end{equation*}
and%
\begin{equation*}
\phi(N) = k \times order(m), \text{ where } k \text{ is a positive integer.}
\end{equation*}
Combining the above results, we get%
\begin{equation*}
m^{\phi(N)} = m^{k \times order(m)} = (m^{order(m)})^k \equiv 1^k \equiv 1 \pmod N\text{.}
\end{equation*}
%
\end{proof}
\end{sectionptx}
%
%
\typeout{************************************************}
\typeout{Section 6 Extended RSA Encryption Scheme}
\typeout{************************************************}
%
\begin{sectionptx}{Section}{Extended RSA Encryption Scheme}{}{Extended RSA Encryption Scheme}{}{}{section-6}
\begin{introduction}{}%
In this section, we define the extended RSA encryption scheme and prove its correctness by demonstrating that the decryption process correctly recovers the original plaintext message.%
\end{introduction}%
\begin{definition}{Definition}{Extended RSA Encryption Scheme.}{def-extended-rsa}%
For any positive integer \(N\), the extended RSA encryption scheme is defined as follows: choose \(e\) and \(d\) such that \(1 \leq e, d \leq \phi(N)\) and \(e \times d \equiv 1 \pmod{\phi(N)}\), then the encryption function is defined as:%
\par
For any \(m\), \(1 \leq m \leq N\), \(\text{Enc}(m) = m^e \pmod{N}\)%
\par
and the decryption function is defined as:%
\par
For any \(c\), \(1 \leq c \leq N\), \(\text{Dec}(c) = c^d \pmod{N}\)%
\end{definition}
Like the classical RSA encryption scheme, the correctness condition for the extended RSA encryption scheme is formulated as:%
\par
%
\begin{equation*}
(m^e)^d \equiv m \pmod{N} \text{ for any integer } m, 1 \leq m \leq N\text{.}
\end{equation*}
%
\par
The main difference between the extended RSA encryption scheme and the classical RSA encryption scheme is that the extended RSA scheme permits a wider choices of \(N\). The main result of this paper shows the set of messages satisfying the correctness condition of the extended RSA encryption scheme is exactly the \(\Phi\text{-set}(N)\). This main result is proved in the next Theorem.%
\begin{theorem}{Theorem}{Correctness of Extended RSA Encryption Scheme.}{}{extended-rsa-correctness}%
For any positive integer \(N\), the correctness condition of the extended RSA encryption scheme holds for all integers in \(\Phi\text{-set}(N)\).%
\par
In other words, for any integer \(m\) in \(\Phi\text{-set}(N)\), we have:%
\begin{equation*}
(m^e)^d \equiv m \pmod{N}\text{.}
\end{equation*}
%
\end{theorem}
\begin{proof}{Proof}{}{extended-rsa-correctness-3}
By definition, for any integer \(m\) in \(\Phi\text{-set}(N)\), we have:%
\begin{equation*}
P = gcd(m,N), Q = N \div P, \text{ and }, gcd(P,Q) = 1.
\end{equation*}
We need to prove for such \(m\) that:%
\begin{equation*}
(m^e)^d \equiv m \pmod{N}
\end{equation*}
%
\par
Since \(e \times d \equiv 1 \pmod{\phi(N)}\), we can express \(e \times d\) as:%
\begin{equation*}
e \times d = k \times \phi(N) + 1
\end{equation*}
for some integer \(k\).%
\par
Now, we can rewrite \((m^e)^d\) as:%
\begin{equation*}
(m^e)^d = m^{k \times \phi(N) + 1}
\end{equation*}
and we can easily verify that:%
\begin{equation*}
m^{k \times \phi(N) + 1} \equiv m \pmod{P}
\end{equation*}
(both sides are divisible by \(P\) due to \(P = gcd(m,N)\)).%
\par
Next, we want to show that \(m^{k \times \phi(N) + 1} \equiv m \pmod{Q}\). Since \(gcd(P,Q) = 1\), \(\phi(N) = \phi(P) \cdot \phi(Q)\) holds and we can rewrite \(m^{k \times \phi(N) + 1}\) as:%
\begin{equation*}
m^{k \times \phi(P) \cdot \phi(Q) + 1} = m^{k \times \phi(P) \cdot \phi(Q)} \cdot m\text{.}
\end{equation*}
%
\par
Since \(gcd(P,Q) = 1\), which implies \(gcd(m,Q) = 1\) (because \(m \in \Phi\text{-set}(N)\)), we can apply Euler's theorem, which states that if \(gcd(m,Q) = 1\), then:%
\begin{equation*}
m^{\phi(Q)} \equiv 1 \pmod{Q}\text{.}
\end{equation*}
From this we can conclude that:%
\begin{equation*}
m^{k \times \phi(P) \cdot \phi(Q) + 1} \equiv m \pmod{Q}\text{.}
\end{equation*}
%
\par
Summarize what we have proved:%
\begin{equation*}
(m^e)^d \equiv m \pmod{P}
\end{equation*}
and%
\begin{equation*}
(m^e)^d \equiv m \pmod{Q}\text{.}
\end{equation*}
Since \(gcd(P,Q) = 1\) and \(N = P \times Q\), by \hyperref[congruence-theorem]{Theorem~{\xreffont\ref{congruence-theorem}}, p.\,\pageref{congruence-theorem}}, we have:%
\begin{equation*}
(m^e)^d \equiv m \pmod{N}\text{.}
\end{equation*}
This completes the proof of the theorem.%
\end{proof}
\begin{corollary}{Corollary}{.}{}{corollary-extended-rsa}%
If \(N\) is a positive integer and%
\begin{equation*}
N = p_1 p_2 \cdots p_k
\end{equation*}
where%
\begin{equation*}
p_1, p_2, \ldots, p_k
\end{equation*}
are distinct prime numbers, then the correctness condition of the RSA encryption scheme holds for all \(m\), \(1 \leq m \leq N\).%
\end{corollary}
This corollary follows directly from \hyperref[extended-rsa-correctness]{Theorem~{\xreffont\ref{extended-rsa-correctness}}, p.\,\pageref{extended-rsa-correctness}} and \hyperref[Phi-Set-Theorem]{Theorem~{\xreffont\ref{Phi-Set-Theorem}}, p.\,\pageref{Phi-Set-Theorem}}.%
\par
In the classical RSA algorithm, where \(N = p \times q\) with \(p\) and \(q\) distinct primes, this corollary establishes the correctness of the classical RSA encryption scheme.%
\begin{example}{Example}{Example of classical RSA Encryption Scheme.}{section-6-11}%
Consider \(N = 10 = 2 \times 5\). We can easily compute:%
\begin{equation*}
\phi\text{-set}(N) = \{1,3,7,9\}
\end{equation*}
and%
\begin{equation*}
\Phi\text{-set}(N) = \{1,2,3,4,5,6,7,8,9,10\}\text{.}
\end{equation*}
If we choose \(e = 3\) and \(d = 7\) (since \(3 \times 7 \equiv 1 \pmod{4}\) and \(\phi(10) = 4\)) as the keys for the RSA encryption scheme, then from the corollary, it follows the correctness condition holds for all \(m\), \(1 \leq m \leq 10\).%
\end{example}
\begin{example}{Example}{Example of extended RSA Encryption Scheme.}{section-6-12}%
Consider \(N = 20 = 2^2 \times 5 \). We can easily compute:%
\begin{equation*}
\phi\text{-set}(N) = \{1,3,7,9,11,13,17,19\}
\end{equation*}
and%
\begin{equation*}
\Phi\text{-set}(N) = \{1,3,4,5,7,8,9,11,12,13,15,16,17,19,20\}\text{.}
\end{equation*}
If we choose \(e = 3\) and \(d = 3\) (since \(3 \times 3 \equiv 1 \pmod{8}\) and \(\phi(20) = 8\)), as the keys for the RSA encryption scheme, then we can verify for any \(m\) in \(\Phi\text{-set}(20)\) that the correctness condition holds. However, \textbraceleft{}2,6,10,14,18\textbraceright{} are not in \(\Phi\text{-set}(20)\) and we can verify that all of those integers do not satisfy the correctness condition. In fact,%
\begin{equation*}
(2^3)^3 \equiv 12 \pmod{20}\text{,}
\end{equation*}
%
\begin{equation*}
(6^3)^3 \equiv 16 \pmod{20}\text{,}
\end{equation*}
%
\begin{equation*}
(10^3)^3 \equiv 0 \pmod{20}\text{,}
\end{equation*}
%
\begin{equation*}
(14^3)^3 \equiv 4 \pmod{20}\text{,}
\end{equation*}
%
\begin{equation*}
(18^3)^3 \equiv 8 \pmod{20}\text{.}
\end{equation*}
None of these results equals the original plaintext message. This example demonstrates that, for the given choices of \(N, e, \text{ and } d\), the messages for which the extended RSA encryption scheme is correct are exactly those in \(\Phi\text{-set}(N)\).%
\end{example}
\end{sectionptx}
%
%
\typeout{************************************************}
\typeout{Section 7 Future Work}
\typeout{************************************************}
%
\begin{sectionptx}{Section}{Future Work}{}{Future Work}{}{}{section-7}
\begin{introduction}{}%
In this section, we will discuss future work related to the generalized RSA encryption scheme.%
\end{introduction}%
%
%
\typeout{************************************************}
\typeout{Subsection 7.1 Necessary Condition of the Correctness \(\Phi\text{-set}(N)\)}
\typeout{************************************************}
%
\begin{subsectionptx}{Subsection}{Necessary Condition of the Correctness \(\Phi\text{-set}(N)\)}{}{Necessary Condition of the Correctness \(\Phi\text{-set}(N)\)}{}{}{subsection-7-1}
We have proved that all messages in \(\Phi\text{-set}(N)\) satisfy the correctness condition of the generalized RSA encryption scheme. However, we did not prove that it is necessary that messages must be in \(\Phi\text{-set}(N)\) to satisfy the correctness condition. Example 6.5 in section 6 demonstrates that the necessary condition for the correctness of the generalized RSA encryption scheme is true. Moreover, we have performed many computational experiments for various values of \(N\) and \(e\) (except \(e = 1\)) to verify that the necessary condition is true for all tested cases. Therefore, it is reasonable to conjecture that the necessary condition for the correctness of the generalized RSA encryption scheme is true. However, a formal proof of this conjecture is still missing and needs to be done in future work.%
\end{subsectionptx}
%
%
\typeout{************************************************}
\typeout{Subsection 7.2 Efficient Computing of \(\Phi(N)\)}
\typeout{************************************************}
%
\begin{subsectionptx}{Subsection}{Efficient Computing of \(\Phi(N)\)}{}{Efficient Computing of \(\Phi(N)\)}{}{}{subsection-7-2}
The computation of \(\Phi(N)\) is another interesting problem. However, the brute-force method for computing \(\Phi(N)\) is not efficient enough for large values of \(N\). In future work, we will explore more efficient algorithms for computing \(\Phi(N)\).%
\end{subsectionptx}
\end{sectionptx}
%
%
\typeout{************************************************}
\typeout{References 8 References}
\typeout{************************************************}
%
\begin{references-section}{References}{References}{}{References}{}{}{references}
%% If this is a top-level references
%%   you can replace the "referencelist" environment
%%   with a standard "thebibliography" environment
\begin{referencelist}
\bibitem[1]{RSA78}\label{RSA78}{}\hypertarget{RSA78}{}Ronald L Rivest, Adi Shamir, and Leonard Adleman, ``A method for obtaining digital signatures and public-key cryptosystems'', in \pubtitle{Communications of the ACM,}, 21(2):120–126, 1978.
\end{referencelist}
\end{references-section}
\end{document}
